\section*{A}
\label{sec:taska}

We ask how likely it is that the disease dies out instead of growing into a big
outbreak. In a closed SIR model every run ends with $I=0$ sooner or later,
because in the end there is no one left to infect. So ``disappear'' cannot mean
``$I$ reaches $0$'' -- that always happens. It means a minor outbreak:
the disease fizzles out early and only a few people are ever infected, instead
of a major epidemic that infects a large part of the population. We
measure the outbreak by the final size $R_\infty$, the total number of people
ever infected, and call a run ``disappeared'' if $R_\infty$ stays below a small
threshold (we use $1\%$ of $N$, i.e.\ a cutoff of $10$ for the $N=1000$ runs
below).

A simple formula predicts this probability. At the very start almost everyone is
susceptible ($S\approx N$), so one infected person makes new infections at rate
$\beta$ and recovers at rate $\gamma$, and before recovering it infects on
average $R_0=\beta/\gamma$ others (\autoref{eq:R0}). While $I$ is still small,
the spread behaves like a simple branching process: each infected person starts
its own chain, and the chains do not yet bump into each other.

Let $q$ be the probability that one such chain eventually dies out. We find $q$
by conditioning on the first event of a single infected person. That
person has two exponential clocks running at once: a recovery clock of rate
$\gamma$ and an infection clock of rate $\tfrac{\beta}{N}S\approx\beta$ (since
$S\approx N$ early on). Whichever rings first is the next event, and for two
competing exponential clocks of rates $a$ and $b$ the ``$a$'' clock wins with
probability $a/(a+b)$ -- the same competing-clocks result we used for the event
selection in \autoref{eq:PEA}. Hence the person recovers first with probability
$\gamma/(\beta+\gamma)$ and infects first with probability $\beta/(\beta+\gamma)$.
\begin{itemize}
  \item If it recovers first, nobody was infected and the chain stops --
        it has died out for certain.
  \item If it infects first, there are now two infectives. Each one
        starts a fresh chain that is statistically identical to the original, so
        each dies out with the same probability $q$; being independent, they
        both die out with probability $q\cdot q=q^2$.
\end{itemize}
Adding the two disjoint cases (and noting $q$ reappears because after an
infection we face the same problem again) gives a self-consistency equation
\begin{equation}
  q = \frac{\gamma}{\beta+\gamma} + \frac{\beta}{\beta+\gamma}\,q^2,
  \label{eq:qfix}
\end{equation}
i.e.\ $\beta q^2 - (\beta+\gamma)q + \gamma = 0$, which factors as
$(q-1)(\beta q-\gamma)=0$ with roots $q=1$ and $q=\gamma/\beta=1/R_0$. The true
extinction probability of a branching process is the smallest root in
$[0,1]$, so
\begin{equation}
  q =
  \begin{cases}
    1/R_0, & R_0 > 1,\\
    1,     & R_0 \le 1.
  \end{cases}
  \label{eq:q}
\end{equation}
The $I_0$ starting infectives spread independently, so the disease disappears
only if all of their chains die out:
\begin{equation}
  P(\text{disease disappears}) = q^{\,I_0} = \left(\frac{1}{R_0}\right)^{I_0}
  \quad (R_0 > 1).
  \label{eq:pdisappear}
\end{equation}
This matches the qualitative picture in the \nameref{sec:ode} section: for $R_0 \le 1$ the
disease always dies out ($P=1$), and the more infectious it is (larger $R_0$),
the less likely it is to disappear.

To test this, we first fix $N=1000$, $R_0=3$ and $I_0=1$ and run $2000$
independent simulations, estimating $P(\text{disease disappears})$ as the fraction of
minor outbreaks with a $95\%$ confidence interval from the replications:
\[
  \hat{P}(\text{disease disappears}) = 0.335 \pm 0.021,
  \qquad
  \text{theory } (1/R_0)^{I_0} = 0.333 .
\]
The two agree. \autoref{fig:final_size} also shows why the question makes sense:
the final size $R_\infty$ is strongly bimodal -- a spike of minor
outbreaks near zero and a separate bump of major epidemics near
$R_\infty\approx 940$, with an empty gap between them. Because of this gap the
exact threshold does not matter.

\begin{figure}[H]
  \centering
  \includegraphics[width=0.8\linewidth]{part1//media/final_size_distribution.png}
  \caption{Final size $R_\infty$ over $2000$ runs ($N=1000$, $R_0=3$, $I_0=1$).
  The distribution is bimodal: minor outbreaks (left) and major epidemics
  (right). Both panels use the same bin width.}
  \label{fig:final_size}
\end{figure}

We then sweep the parameters to check \autoref{eq:pdisappear} more broadly
($1000$ runs per point), shown in \autoref{fig:extinction_sweep}. Sweeping $R_0$
(left, $I_0=1$), $P(\text{disease disappears})$ drops from about $1$ towards $0$ as the
disease becomes more infectious and follows the $(1/R_0)^{I_0}$ curve. Sweeping
$I_0$ (right, $R_0=3$), it falls geometrically and is a straight line on the log
axis, which confirms the exponent $I_0$ in \autoref{eq:pdisappear}.

\begin{figure}[H]
  \centering
  \includegraphics[width=0.8\linewidth]{part1//media/extinction_sweep.png}
  \caption{$P(\text{disease disappears})$ versus $R_0$ (left) and versus $I_0$ (right,
  log scale), with $95\%$ confidence intervals. Dashed line is the theory
  $(1/R_0)^{I_0}$.}
  \label{fig:extinction_sweep}
\end{figure}

Two small mismatches are worth noting. Near $R_0=1$ the simulation sits a little
below the theory (e.g.\ $R_0=1.2$: $0.76$ vs.\ $0.83$), because the branching
formula assumes an endless supply of susceptibles; close to $R_0=1$ outbreaks
are slow and the finite population makes them die out more often than the
formula predicts. For $R_0<1$ the estimate is not exactly $1$ either, because
with a finite population a few runs still cross the small threshold by chance.
Both effects are finite-population artefacts and should vanish as $N\to\infty$,
where the endless-susceptibles assumption becomes exact. \autoref{fig:finite_n}
confirms this: fixing each of the two awkward cases ($R_0=0.8$ and $R_0=1.2$,
$I_0=1$) and sweeping $N$, the estimate climbs monotonically onto the theory line
as $N$ grows (for $R_0=1.2$, from $0.58$ at $N=250$ to $0.83$ at $N=8000$,
matching $1/R_0=0.833$). This is also the mechanism behind question~(d): the
larger the population, the better the susceptible pool matches the
infinite-population assumption that the deterministic and branching descriptions
rely on.

\begin{figure}[H]
  \centering
  \includegraphics[width=0.5\linewidth]{part1//media/extinction_finite_n.png}
  \caption{$P(\text{disease disappears})$ versus population size $N$ (log axis, $1000$ runs
  per point, $I_0=1$) for the two cases where the branching formula is least
  accurate. Dashed lines are the population-independent theory $(1/R_0)^{I_0}$.
  Both curves converge onto the theory as $N$ grows, confirming that the
  mismatches in \autoref{fig:extinction_sweep} are finite-population effects.}
  \label{fig:finite_n}
\end{figure}
\section*{B}
\label{sec:taskb}

Here we ask a simple question: can the disease come back again and again in
repeated waves, instead of just sweeping through once? In the plain SIR model the
answer is no. The reason is easy to see: the number of people who can still catch
the disease, $S$, only ever goes down. People move from $S$ to $I$ to $R$ and
never come back. So once everyone has been infected and recovered, there is no
one left to infect, and the outbreak is a single wave that ends at $I=0$ and stays
there (\autoref{fig:sir_wave}). Nothing refills $S$, so a second wave can never
start.

\begin{figure}[H]
  \centering
  \includegraphics[width=0.5\linewidth]{part1//media/sir_single_run.png}
  \caption{One stochastic SIR run ($N=1000$, $R_0=3$). It is a single wave: $S$
  only falls, $I$ goes up once and then back to $0$. No cycles are possible
  because no one ever becomes susceptible again.}
  \label{fig:sir_wave}
\end{figure}

To get repeated waves, we need a way to put people back into $S$. The easiest one
is fading immunity: after a while, recovered people lose their protection and can
catch the disease again. Now, we are not doctors or anything, but this seems to be
roughly what happens with diseases like the flu or COVID -- the protection you get
from being sick does not last forever. This just adds one new move, $R\to S$,
happening at rate $\omega$. The model is now called SIRS. It is the same SIR chain as before, but
with one extra arrow that loops $R$ back to $S$:

\begin{center}
\begin{tikzpicture}[
    node distance=22mm and 26mm,
    comp/.style={circle, draw, minimum size=11mm, font=\large},
    every edge/.style={draw, -{Latex[length=2.5mm]}}
  ]
  \node[comp] (S) {$S$};
  \node[comp, right=of S] (I) {$I$};
  \node[comp, right=of I] (R) {$R$};

  \draw (S) edge node[above]{$\tfrac{\beta}{N}SI$} (I);
  \draw (I) edge node[above]{$\gamma I$} (R);
  \draw (R) edge[bend left=50] node[below]{$\omega R$ (waning)} (S);
\end{tikzpicture}
\end{center}

\noindent Written as equations it is
\begin{equation}
  \left\{
  \begin{aligned}
    \frac{dS}{dt} &= -\frac{\beta\,S\,I}{N} + \omega R,\\[4pt]
    \frac{dI}{dt} &= \frac{\beta\,S\,I}{N} - \gamma I,\\[4pt]
    \frac{dR}{dt} &= \gamma I - \omega R.
  \end{aligned}
  \right.
  \label{eq:sirs_ode}
\end{equation}
which is the same as the SIR equations (\autoref{eq:ode}) plus the one extra term
$\omega R$ that sends recovered people back to $S$. Because people keep coming
back to $S$, the disease no longer has to disappear. If $R_0>1$
(\autoref{eq:R0}), the disease can stay around forever at a steady level, where
the number of infected people sits at
\begin{equation}
  S^\star = \frac{N}{R_0},
  \qquad
  I^\star = N\Bigl(1-\tfrac{1}{R_0}\Bigr)\frac{\omega}{\omega+\gamma}.
  \label{eq:endemic}
\end{equation}
We call $I^\star$ the steady (endemic) level.

What does this steady level look like as the system settles down? If we just
solve the equations (\autoref{eq:sirs_ode}), the number of infected people
wobbles up and down a few times, with each wobble smaller than the last, until it
flattens out at $I^\star$. So the plain equations do make waves, but the waves
shrink and die away. More precisely, it can be shown that linearising
\autoref{eq:sirs_ode} about the endemic equilibrium \eqref{eq:endemic} yields a
Jacobian whose eigenvalues have negative real part, so the equilibrium is stable.
For our parameters (in particular a small waning rate $\omega$) these eigenvalues
form a complex-conjugate pair, making the equilibrium a stable
focus,\footnote{This linear stability classification -- real eigenvalues
(node, overdamped) versus a complex-conjugate pair (focus, damped oscillation) --
is exactly the kind of analysis some of us saw in 01617 Introduction to Dynamical
Systems.} so $I(t)$
returns to $I^\star$ through damped oscillations whose amplitude decays
exponentially -- for a large enough $\omega$ the eigenvalues are instead real and
the approach is overdamped, with no oscillation at all. Either way the
deterministic model predicts only a transient, not a sustained cycle -- after it
dies out, $I(t)\to I^\star$ constant.

The simulation behaves differently, and this is the interesting part.
\autoref{fig:sirs_cycles} shows one SIRS simulation ($N=20\,000$, $R_0=5$,
$\gamma=0.1$, $\omega=0.002$) started near the steady level, drawn on top of the
smooth equation curve that starts from the same place. At the start, when the
system is still far from the steady level, the simulation closely follows the
smooth curve. But once it gets near the steady level, something changes. In a real
(finite) population, people get infected and recover one at a time, at random
moments. These little random bumps keep nudging the system, so instead of
flattening out, the simulated number of infected people keeps wobbling up and down
forever. The smooth equation has no randomness, so it flattens; the simulation
never does.

The two curves wobble about as often, but not at the same times, so after the
first wave they stop lining up -- the random bumps slowly shift the timing of the
simulated waves. The picture on the right shows the same thing in another way: the
smooth curve spirals inward and stops at the steady point, while the noisy
simulation keeps circling around it and never settles.

\begin{figure}[H]
  \centering
  \includegraphics[width=0.8\linewidth]{part1//media/sirs_cycles.png}
  \caption{SIRS near the steady level ($N=20\,000$, $R_0=5$, $\gamma=0.1$,
  $\omega=0.002$). Left: infected over time, simulation (red) vs. the smooth
  equation (black). The equation flattens out at $I^\star$, the simulation keeps
  making waves. Right: the same runs drawn as $S$ against $I$ -- the equation
  spirals into the steady point, the simulation keeps circling it. Both start from
  the same state.}
  \label{fig:sirs_cycles}
\end{figure}

So yes, the disease can show cyclical behaviour -- but only in the simulation, not
in the plain equations. Just like the disease dying out by chance in
\nameref{sec:taska}, this is something the smooth equations miss completely. The
equations let the waves fade away -- the real, finite population keeps them going,
because real people come in whole numbers and get sick at random times. This is
the same lesson as in question~(a): when the population is finite,
randomness can change the whole behaviour, not just the small details. One thing
to keep in mind: in a small population, a random dip can take $I$ all the way to
$0$, and then the disease is gone for good (nothing brings it back). So the waves
only keep going if the population is big enough that $I^\star$ stays well above
zero, which is why we use a large $N=20\,000$ here.